% !TEX root = main.tex
% ============================================================================
% CHAPTER 1: INTRODUCTION AND RESEARCH SPECIFICATIONS
% ============================================================================

\chapter{کلیات تحقیق}
\section{بيان مسئله}

در عصر دیجیتال، تبادل اطلاعات به‌ویژه در قالب تصاویر، به بخشی جدایی‌ناپذیر از ارتباطات، تجارت، سیستم‌های نظارتی، و رسانه‌های اجتماعی تبدیل شده است. در این میان، تصاویر دیجیتال به دلیل نقش کلیدی در انتقال داده‌های بصری، به یکی از مهم‌ترین انواع اطلاعات چندرسانه‌ای بدل شده‌اند که حجم بالایی از داده را در خود جای می‌دهند و کاربرد گسترده‌ای در شبکه‌های ارتباطی، سامانه‌های ذخیره‌سازی و تحلیل داده دارند \cite{ref1}.

با گسترش دسترسی به فضای مجازی و تهدیدهای فزاینده در حوزه امنیت سایبری، حفاظت از حریم خصوصی اطلاعات تصویری به یکی از دغدغه‌های اصلی پژوهشگران و توسعه‌دهندگان امنیت اطلاعات تبدیل شده است. برخلاف داده‌های متنی که رمزنگاری آن‌ها با الگوریتم‌های استانداردی نظیر استاندارد رمزنگاری پیشرفته\footnote{\lr{Advanced Encryption Standard (AES)}} و الگوریتم ریوسـت-شامیر-آدلمن\footnote{\lr{Rivest-Shamir-Adleman (RSA)}} به‌خوبی قابل انجام است، تصاویر دیجیتال به دلیل ویژگی‌هایی همچون حجم زیاد، ساختار دوبعدی، همبستگی آماری بالا میان پیکسل‌های مجاور، و حساسیت بالا به تغییرات، به الگوریتم‌های رمزنگاری تخصصی‌تر و پیچیده‌تری نیاز دارند \cite{ref2, ref3}.

مهم‌ترین چالش‌ها در رمزنگاری تصاویر دیجیتال را می‌توان به شرح زیر برشمرد:

\begin{itemize}
    \item حفظ ویژگی‌های آماری تصویر پس از رمزنگاری
    \item ایجاد مقاومت در برابر انواع حملات نظیر تحلیل آماری، حملات تمایزی و جستجوی فراگیر
    \item کاهش بار محاسباتی جهت کاربرد در سامانه‌های بلادرنگ
    \item پیشگیری از تخریب داده‌های تصویری طی فرآیند رمزنگاری و رمزگشایی
\end{itemize}

در پاسخ به این چالش‌ها، بهره‌گیری از سیستم‌های آشوبی به‌عنوان راهکاری مؤثر مورد توجه قرار گرفته است. سیستم‌های آشوبی با رفتار غیرخطی، حساسیت شدید به شرایط اولیه، و توانایی تولید دنباله‌های شبه‌تصادفی، امکان طراحی الگوریتم‌هایی با امنیت بالا و کلیدهای غیرقابل پیش‌بینی را فراهم می‌آورند \cite{ref4}. فرآیندهای جایگشت و پراکندگی پیکسل‌ها با کمک این سیستم‌ها، ساختار آماری تصویر را به‌شدت تغییر داده و مانع بازیابی تصویر اصلی بدون کلید می‌شوند.

با این حال، یکی از نقاط ضعف رایج در الگوریتم‌های مبتنی بر سیستم‌های آشوبی کلاسیک در پیاده‌سازی‌های سخت‌افزاری و نرم‌افزاری، بروز پدیده‌ای به‌نام تخریب آشوب\footnote{\lr{Chaos Degradation}} است \cite{ref5, ref19}. این پدیده ناشی از محدودیت دقت شناور و نمایش دیجیتال اعداد در کامپیوترهاست که باعث می‌شود سیستم‌های آشوبی پیوسته و گسسته پس از طی دوره‌های مشخص، رفتار بی‌نظم خود را از دست داده و به مدارهای دوره‌ای کوتاه\footnote{\lr{Periodic Orbits}} سقوط کنند. این امر موجب کاهش شدید بی‌نظمی، کاهش فضای کلید موثر و در نتیجه آسیب‌پذیری سیستم در برابر حملات تحلیل آماری می‌شود. این اشکال به‌ویژه در نگاشت‌های یک‌بعدی گسسته کلاسیک نظیر نگاشت لجستیک\footnote{\lr{Logistic Map}} و نگاشت تنت\footnote{\lr{Tent Map}} به‌صورت برجسته‌تری مشاهده شده است \cite{ref5}.

سیستم آشوبی به‌طور کلی، نوعی سیستم دینامیکی غیرخطی است که رفتار آن در طول زمان به شدت به شرایط اولیه وابسته است و این وابستگی منجر به تولید پاسخ‌هایی کاملا متفاوت در صورت تغییرات بسیار جزئی در مقدارهای ورودی می‌شود. به بیان ساده‌تر، سیستم آشوبی با وجود پیروی از معادلات مشخص، خروجی‌ای شبه‌تصادفی و غیرقابل پیش‌بینی تولید می‌کند، به‌گونه‌ای که برای ناظر بیرونی همانند نویز یا تصادف به نظر می‌رسد، اما در واقع دارای نظم در بی‌نظمی است.

ویژگی‌های اصلی سیستم‌های آشوبی شامل:

\begin{itemize}
    \item حساسیت شدید به شرایط اولیه
    \item غیرخطی بودن روابط سیستم
    \item رفتار شبه‌تصادفی
    \item وابستگی به پارامترهای سیستمی
\end{itemize}

این ویژگی‌ها باعث می‌شوند که سیستم‌های آشوبی گزینه‌ای بسیار مناسب برای تولید کلیدهای رمزنگاری، الگوریتم‌های پراکندگی و اختفاء داده‌ها باشند. به‌عنوان مثال، نگاشت‌های ساده‌ای مانند لجستیک، تنت، هنون، چن\footnote{\lr{Chen}} و لورنز\footnote{\lr{Lorenz}} به‌طور گسترده در رمزنگاری تصویر به‌کار رفته‌اند. سیستم لورنز یکی از مشهورترین مدل‌های آشوبی سه‌بعدی است که از معادلات دیفرانسیل غیرخطی تشکیل شده و در ابتدا برای مدل‌سازی حرکت سیالات جوی طراحی شد، اما به‌دلیل توانایی بالا در تولید مسیرهای آشوبی، امروزه در حوزه رمزنگاری و ارتباطات امن نیز مورد استفاده قرار گرفته است \cite{ref6}.

از سوی دیگر، پژوهشگران برای افزایش مقاومت در برابر حملات تحلیل همبستگی و جلوگیری از تخریب آشوب، سیستم‌هایی تحت عنوان نگاشت‌های آشوب نمایی را توسعه داده‌اند. این سیستم‌ها، با استفاده از نگاشت‌های ترکیبی و تغییرات پارامتری گسترده، دامنه آشوبی پایدارتری نسبت به نگاشت‌های کلاسیک دارند و در نتیجه، در طراحی الگوریتم‌های رمزنگاری مقاوم بسیار مؤثر واقع شده‌اند \cite{ref4}.

در پژوهش حاضر، با هدف ارتقاء سطح امنیت و کارایی الگوریتم‌های رمزنگاری تصویر، روشی جدید معرفی می‌شود که از ترکیب نه سیستم آشوبی نمایی به‌همراه الگوریتم جایگشت چن بهره می‌برد. سیستم‌های نمایی، مانند \lr{LEL}\footnote{\lr{Logistic-Exponential-Logistic}}، \lr{SEL}\footnote{\lr{Sine-Exponential-Logistic}} و سایر نگاشت‌های آشوب نمائی، به‌عنوان نسل جدید سیستم‌های آشوبی، قادرند رفتارهای پیچیده‌تر و توزیع آماری یکنواخت‌تری نسبت به مدل‌های کلاسیک ایجاد کنند \cite{ref4}.

در این الگوریتم، ابتدا پیکسل‌های تصویر با استفاده از دنباله‌های تولیدشده توسط سیستم چن به‌صورت تصادفی جابه‌جا می‌شوند. سپس، رمزنگاری هر پیکسل با انتخاب تصادفی یکی از ۹ سیستم نمایی صورت می‌گیرد؛ به این ترتیب که مقادیر رنگی پیکسل‌ها با اعداد تولید شده توسط آشوب ترکیب می‌شوند. این رویکرد چندلایه، علاوه بر ارتقاء پیچیدگی، موجب افزایش آنتروپی اطلاعاتی، کاهش همبستگی پیکسل‌ها و افزایش مقاومت در برابر حملات تحلیل آماری می‌شود. از طرفی، با بهره‌گیری از توزیع‌های آماری غیرخطی‌تر و پایدارتر، احتمال تخریب ویژگی‌های آشوبی به‌طور مؤثری کاهش می‌یابد.

در نهایت، هدف از این تحقیق، توسعه یک الگوریتم رمزنگاری مقاوم و کارآمد است که با تکیه بر سیستم‌های آشوبی نوین، توازنی بهینه میان امنیت بالا و سرعت اجرای مناسب برقرار کرده و پاسخگوی نیازهای کاربردی در حوزه‌هایی همچون پزشکی، امنیت نظامی، شبکه‌های اجتماعی و ذخیره‌سازی ابری باشد.

\section{اهمیت و ضرورت تحقیق}

با گسترش فناوری‌های دیجیتال و افزایش وابستگی جوامع به سامانه‌های الکترونیکی، امنیت اطلاعات به یکی از دغدغه‌های اصلی در حوزه فناوری اطلاعات تبدیل شده است. تصاویر دیجیتال، به‌ویژه تصاویر رنگی، به‌دلیل حجم بالای اطلاعات و کاربرد گسترده در حوزه‌هایی مانند ارتباطات، شبکه‌های اجتماعی، نظارت تصویری، پزشکی و نظامی، بیش از گذشته در معرض تهدیدات امنیتی نظیر شنود، جعل و دستکاری قرار دارند \cite{ref10}. اهمیت حفاظت از این داده‌های بصری، ضرورت توسعه روش‌های نوین رمزنگاری را دوچندان کرده است.

یکی از چالش‌های اساسی در رمزنگاری تصاویر رنگی، ساختار پیچیده و چندبعدی آن‌هاست. تصاویر رنگی به‌صورت ترکیبی از سه کانال \lr{RGB}\footnote{\lr{Red-Green-Blue}} تعریف می‌شوند که وابستگی بین کانال‌ها، کار رمزنگاری ایمن و بهینه را دشوار می‌سازد. این وابستگی‌ها می‌توانند توسط مهاجمان برای تحلیل ساختاری تصویر مورد سوءاستفاده قرار گیرند \cite{ref3}.  بنابراین، طراحی الگوریتم‌هایی که بتوانند این وابستگی‌ها را کاهش داده و توزیع داده‌ها را تصادفی‌تر سازند، امری حیاتی به نظر می‌رسد.

مطالعات متعددی در حوزه رمزنگاری آشوبی انجام شده‌اند که نشان می‌دهند سیستم‌های آشوبی به‌دلیل ویژگی‌هایی نظیر حساسیت بالا به شرایط اولیه، رفتار غیرخطی و قابلیت تولید دنباله‌های شبه‌تصادفی پیچیده، ظرفیت بالایی در تولید کلیدهای رمزنگاری غیرقابل پیش‌بینی دارند. با این حال، اغلب روش‌های موجود تنها از یک سیستم آشوبی منفرد استفاده می‌کنند که در برابر حملاتی مانند تحلیل همبستگی و جستجوی فراگیر آسیب‌پذیر هستند \cite{ref12}.  به‌منظور مقابله با این آسیب‌پذیری‌ها، استفاده از ترکیب چند سیستم آشوبی نمایی به‌صورت تصادفی، می‌تواند پیچیدگی رمزنگاری را افزایش داده و مقاومت آن را در برابر انواع حملات ارتقا بخشد.

بر اساس یافته‌های اخیر، سیستم‌های آشوبی نمایی نسبت به سیستم‌های آشوبی کلاسیک مانند نگاشت لجستیک و نگاشت تنت، توزیع احتمال پیچیده‌تری تولید می‌کنند و پایداری بیشتری در برابر تغییرات پارامترها دارند. این ویژگی‌ها آن‌ها را به گزینه‌ای مناسب برای پیاده‌سازی رمزنگاری‌های مقاوم تبدیل کرده است \cite{ref5}. از طرف دیگر، استفاده از سیستم‌هایی نظیر سیستم آشوبی چن، با ویژگی‌هایی همچون سرعت پردازش بالا و سادگی پیاده‌سازی، می‌توانند به افزایش کارایی الگوریتم‌های رمزنگاری کمک کنند.

افزون بر این، بسیاری از کاربردهای حیاتی نظیر انتقال تصاویر پزشکی، داده‌های مالی در بانکداری، یا اطلاعات طبقه‌بندی‌شده در صنایع نظامی، نیازمند الگوریتم‌هایی هستند که هم امنیت بسیار بالایی داشته باشند و هم بتوانند در زمان محدود پردازش انجام دهند؛ این مسئله، لزوم دستیابی به توازنی منطقی میان امنیت و کارایی را پررنگ‌تر می‌کند \cite{ref5}. در همین راستا، پژوهش حاضر با ترکیب ۹ سیستم آشوبی نمایی و استفاده از سیستم آشوبی چن جهت جایگشت، در پی توسعه مدلی نوین است که بتواند این توازن را به‌صورت مؤثر محقق سازد.

در نهایت، رمزنگاری آشوبی به‌عنوان یکی از حوزه‌های نوظهور در علوم رمزنگاری، همچنان نیازمند تحقیقات عمیق‌تری برای بهره‌برداری کامل از ظرفیت‌های خود است. مطالعات نظری و تجربی در این زمینه می‌توانند زمینه‌ساز توسعه مدل‌های امن‌تر و بهینه‌تری برای محافظت از داده‌های حساس در عصر دیجیتال باشند. پژوهش حاضر، با تمرکز بر طراحی و ارزیابی یک روش ترکیبی پیشرفته، می‌کوشد گامی مؤثر در مسیر توسعه دانش رمزنگاری آشوبی بردارد.

\section{اهداف تحقيق}

\subsection{هدف کلی}

طراحی و توسعه‌ی الگوریتمی برای رمزنگاری تصاویر رنگی با بهره‌گیری از ترکیب ۹ سیستم آشوبی نمایی، با هدف ارتقاء امنیت اطلاعات تصویری و دستیابی به سرعت اجرای مناسب و قابل‌قبول برای کاربردهای عملی.

\subsection{اهداف جزئی}

\begin{itemize}
    \item بررسی عملکرد ۹ سیستم آشوبی نمایی در تولید دنباله‌های تصادفی پیچیده و غیرقابل پیش‌بینی
    \item بهبود مقاومت الگوریتم رمزنگاری در برابر حملات آماری و جستجوی فراگیر از طریق ترکیب چندین سیستم آشوبی
    \item افزایش پیچیدگی فرآیند رمزنگاری با استفاده از انتخاب تصادفی سیستم آشوبی برای رمزنگاری هر پیکسل
    \item مقایسه امنیت و کارایی الگوریتم پیشنهادی با سایر روش‌های رمزنگاری مبتنی بر آشوب و ارزیابی میزان برتری آن
\end{itemize}

\section{فرضیه‌های تحقیق}

\textbf{فرضیه ۱:} بهره‌گیری از ترکیب ۹ سیستم آشوبی نمایی در فرآیند رمزنگاری سبب افزایش مقاومت الگوریتم در برابر حملات آماری و تحلیل تفاضلی\footnote{\lr{Differential Cryptanalysis}: روشی تحلیل رمز که در آن با بررسی تفاوت میان چند متن رمزشده حاصل از متن‌های اصلی مشابه (با تغییر جزئی)، الگوهای آماری غیرتصادفی برای استخراج کلید یا اطلاعات ساختاری الگوریتم شناسایی می‌شود.} می‌شود.

\textbf{فرضیه ۲:} استفاده از سیستم آشوبی سه‌بعدی چن برای جایگشت، همبستگی میان پیکسل‌های مجاور تصویر را به صفر نزدیک می‌کند.

\textbf{فرضیه ۳:} انتخاب تصادفی و مستقل یکی از ۹ نگاشت آشوبی نمایی برای هر پیکسل، آنتروپی اطلاعاتی تصویر رمزشده را به مقدار ایده‌آل (۸ بیت) نزدیک می‌کند.

\textbf{فرضیه ۴:} ساختار پنج‌مرحله‌ای پیشنهادی با وجود چندلایه بودن، سرعت اجرای مناسب و قابل‌قبولی برای کاربردهای عملی دارد.

\textbf{فرضیه ۵:} اعمال لایه انتشار دوم بر پایه نگاشت لجستیک، حساسیت تصویر رمزشده را نسبت به تغییرات جزئی در کلید یا تصویر اصلی افزایش می‌دهد.

\section{مفاهیم کلیدی تحقیق}

رمزنگاری:
فرایندی ریاضیاتی است که طی آن، داده‌های اصلی (تصویر، متن، صدا و ...) به فرم غیرقابل فهم و ایمن تبدیل می‌شوند، به‌گونه‌ای که تنها افراد مجاز با استفاده از کلید یا الگوریتم رمزگشایی می‌توانند به اطلاعات اصلی دسترسی پیدا کنند. هدف اصلی رمزنگاری، حفظ محرمانگی، صحت و تمامیت اطلاعات در برابر دسترسی یا دستکاری غیرمجاز است.

رمزنگاری تصویر :

زیرشاخه‌ای از رمزنگاری است که تمرکز آن بر رمزنگاری داده‌های تصویری به‌ویژه تصاویر دیجیتال رنگی یا خاکستری است. رمزنگاری تصویر باید با در نظر گرفتن ویژگی‌های خاص تصویر نظیر حجم بالا، وابستگی بین پیکسل‌ها، ساختار دو یا سه‌بعدی و حساسیت به تغییرات طراحی شود.

تصاویر رنگی:

تصاویری هستند که هر پیکسل آن‌ها از ترکیب سه کانال اصلی رنگی قرمز (\lr{R})، سبز (\lr{G}) و آبی (\lr{B}) تشکیل شده است. این کانال‌ها مقادیر عددی خاصی دارند (معمولا در بازه ۰ تا ۲۵۵) که در مجموع رنگ نهایی پیکسل را مشخص می‌کنند. رمزنگاری این نوع تصاویر نیازمند رویکردهای پیچیده‌تری نسبت به تصاویر خاکستری است، زیرا اطلاعات بین کانال‌ها وابسته و همبسته هستند.

سیستم آشوبی:

سیستم‌های دینامیکی غیرخطی هستند که رفتار آن‌ها در طول زمان به شدت به شرایط اولیه وابسته است. این سیستم‌ها با وجود آن‌که معادلات دقیق و قطعی دارند، نتایجی شبه‌تصادفی، غیرقابل پیش‌بینی و پیچیده تولید می‌کنند. ویژگی‌هایی مانند حساسیت به مقدار اولیه، غیرخطی‌بودن و قابلیت تولید دنباله‌های پیچیده، آن‌ها را برای تولید کلیدها و ایجاد اختفاء در رمزنگاری بسیار مناسب ساخته است.

نگاشت آشوب نمایی:

دسته‌ای پیشرفته از سیستم‌های آشوبی هستند که با ترکیب نگاشت‌های ریاضی ساده مانند لجستیک، سینوسی و تنت و افزودن مولفه‌های نمایی، رفتارهای غیرخطی و پویاتری را ایجاد می‌کنند. این نگاشت‌ها دارای دامنه آشوب وسیع‌تر، پایداری بیشتر نسبت به خطا و قابلیت تولید دنباله‌های رمزنگاری غیرقابل پیش‌بینی هستند. نمونه‌هایی از این سیستم‌ها عبارتند از \lr{LEL}، \lr{SEL}، \lr{LET}\footnote{\lr{Logistic-Exponential-Tent}}  و ...

جایگشت :
یکی از مراحل اصلی در رمزنگاری تصویر است که طی آن موقعیت پیکسل‌های تصویر به‌صورت غیرقابل پیش‌بینی تغییر داده می‌شود، بدون آن‌که مقادیر رنگی آن‌ها تغییر یابد. هدف این مرحله، از بین بردن ساختار و همبستگی مکانی بین پیکسل‌ها است تا تحلیل آماری برای مهاجم دشوار شود.

انتشار:
مرحله‌ای از رمزنگاری که در آن مقادیر رنگی پیکسل‌ها تغییر داده می‌شوند. به عبارتی، اطلاعات تصویری با مقادیر تصادفی ترکیب می‌شود تا رابطه‌ای میان داده‌های ورودی و خروجی باقی نماند. این مرحله باعث افزایش مقاومت در برابر تحلیل تفاضلی می‌شود.

آنتروپی شانون :

مقیاسی از میزان تصادفی‌بودن یا بی‌نظمی در توزیع داده‌هاست. در رمزنگاری تصویر، هر چه مقدار آنتروپی به مقدار ایده‌آل (برای تصاویر ۸ بیتی، مقدار ۸) نزدیک‌تر باشد، نشان‌دهنده آن است که الگوریتم رمزنگاری به‌خوبی توانسته تصویر را غیرقابل تحلیل کند.

نرخ تغییر تعداد پیکسل‌ها (\lr{NPCR})\footnote{\lr{Number of Pixel Change Rate}} و میانگین شدت تغییرات یکپارچه (\lr{UACI})\footnote{\lr{Unified Average Changing Intensity}}:

دو معیار استاندارد و بنیادی برای ارزیابی میزان حساسیت الگوریتم رمزنگاری نسبت به تغییرات جزئی (حتی در حد یک پیکسل یا یک بیت) در کلید یا تصویر اصلی هستند.

\textbf{نرخ تغییر تعداد پیکسل‌ها (\lr{NPCR}):} نشان‌دهنده درصد تعداد پیکسل‌هایی است که در اثر تغییر جزئی در ورودی، در تصویر رمزنگاری‌شده دچار تغییر می‌شوند.

\textbf{میانگین شدت تغییرات یکپارچه (\lr{UACI}):} نشان‌دهنده میانگین شدت و تفاوت تراز خاکستری/رنگی میان دو تصویر رمزنگاری‌شده در اثر تغییر جزئی در ورودی است.

\section{سهم علمی پژوهش}

در این پایان‌نامه، یک چارچوب نوین برای رمزنگاری تصاویر رنگی مبتنی بر سیستم‌های آشوبی نمایی ارائه شده است که با هدف افزایش امنیت، پیچیدگی دینامیکی و معکوس‌پذیری کامل طراحی گردیده است. مهم‌ترین سهم‌های علمی این پژوهش به شرح زیر می‌باشند:

\begin{itemize}
    \item ارائه یک الگوریتم رمزنگاری تصویر مبتنی بر مدل آشوبی نمایی که با ترکیب غیرخطی نگاشت‌های پایه، دنباله‌های آشوبی با رفتار پیچیده‌تر و حساسیت بالاتر نسبت به شرایط اولیه تولید می‌کند
    \item استفاده هم‌زمان از چند سیستم آشوبی نمایی (۹ نگاشت متفاوت) و انتخاب پویاى آن‌ها در سطح پیکسل، که منجر به افزایش فضای کلید، کاهش الگوپذیری و ارتقای مقاومت الگوریتم در برابر حملات آماری و دیفرانسیلی می‌شود
    \item طراحی یک سازوکار انتشار اصلاح‌شده و کاملا معکوس‌پذیر که در آن دنباله آشوبی صرفا به کلید و اندیس پیکسل وابسته بوده و از وابستگی مستقیم به مقدار پیکسل جلوگیری شده است. این ویژگی تضمین می‌کند که تصویر رمزگشایی‌شده به‌صورت صددرصد با تصویر اصلی منطبق باشد
    \item بهره‌گیری از سیستم آشوبی پیوسته چن برای تولید جریان کلید، که در مقایسه با نگاشت‌های گسسته، کیفیت آشوب، حساسیت به شرایط اولیه و امنیت کلید را به‌طور قابل‌توجهی افزایش می‌دهد
    \item انجام مجموعه‌ای جامع از تحلیل‌های امنیتی شامل آنتروپی، همبستگی پیکسل‌های مجاور، نرخ تغییر تعداد پیکسل‌ها، میانگین شدت تغییرات یکپارچه و تحلیل هیستوگرام، که نشان‌دهنده عملکرد مناسب الگوریتم پیشنهادی در مقایسه با روش‌های موجود می‌باشد
    \item پیاده‌سازی کامل الگوریتم و ارزیابی تجربی آن بر روی تصاویر رنگی استاندارد با ابعاد مختلف، که کارایی عملی و قابلیت کاربرد الگوریتم در سناریوهای واقعی را تأیید می‌کند
\end{itemize}

\section{ساختار پایان‌نامه}

ساختار این پایان‌نامه به‌صورت زیر سازمان‌دهی شده است:

در فصل دوم، مفاهیم پایه و پیش‌نیازهای نظری مرتبط با رمزنگاری تصاویر، سیستم‌های آشوبی، ویژگی‌های نگاشت‌های آشوبی و معیارهای ارزیابی امنیت مرور می‌شود. همچنین برخی از مهم‌ترین پژوهش‌های انجام‌شده در این حوزه مورد بررسی قرار می‌گیرند.

فصل سوم به معرفی الگوریتم پیشنهادی اختصاص دارد. در این فصل، معماری کلی روش ارائه‌شده، فرآیند تولید کلید، مرحله جایگشت و مرحله انتشار به‌طور دقیق تشریح شده و روند رمزنگاری و رمزگشایی به‌صورت گام‌به‌گام توضیح داده می‌شود.

در فصل چهارم، پیاده‌سازی عملی الگوریتم و نتایج آزمایش‌های تجربی ارائه می‌گردد. در این فصل، عملکرد الگوریتم پیشنهادی از نظر معیارهای امنیتی و آماری تحلیل شده و نتایج به‌صورت جداول و نمودارهای تحلیلی گزارش می‌شوند.

در فصل پنجم، نتایج به‌دست‌آمده مورد بحث و بررسی قرار گرفته و عملکرد روش پیشنهادی با برخی از روش‌های مطرح در ادبیات پژوهش مقایسه می‌شود. در پایان این فصل، جمع‌بندی نهایی پژوهش ارائه شده و پیشنهادهایی برای تحقیقات آینده مطرح می‌گردد.